\documentclass[12pt]{book}
\usepackage[utf8]{inputenc}
\usepackage[T1]{fontenc}
\usepackage{tocbibind}
\usepackage{mathptmx}
\usepackage{geometry}
\usepackage{mathtools}
\usepackage[english]{babel}
\usepackage{graphicx}
\usepackage{subcaption}
\usepackage{stackengine}
\usepackage[os=win]{menukeys}
\usepackage{hyperref}
\usepackage{xcolor}
\usepackage{color}
\usepackage{tikz}
\usepackage[yyyymmdd,hhmmss]{datetime}
\usepackage{etoolbox}
\usepackage[inline]{enumitem}
\usepackage{listings}
\usepackage{booktabs}
\usepackage[os=win]{menukeys}
\usepackage{minted}

\newcommand{\WindowsLogo}{\raisebox{-0.1em}{\includegraphics[height=0.8em]{images/logo/Windows_3_logo_simplified}}}
%\newcommand{\PowerLogo}{\raisebox{-0.1em}{\includegraphics[height=0.8em]{images/logo/power}}}
\newcommand{\WinKey}{\keys{\WindowsLogo}}
\newcommand{\PowerKey}{\keys{\PowerLogo}}

%%%%% Mengganti label "Contents" ke "Daftar Isi" %%%%%
\addto\captionsenglish{\renewcommand{\contentsname}{Daftar Isi}}

%%%%% Mengganti label "Chapter" ke "Bab" %%%%%
\addto\captionsenglish{\renewcommand{\chaptername}{Bab}}

%%%%% Mengganti label "Figure" ke "Gambar" %%%%%
\addto\captionsenglish{\renewcommand{\figurename}{Gambar}}

%%%%% Mengganti label "List of Figures" ke "Daftar Gambar" %%%%%
\addto\captionsenglish{\renewcommand{\listfigurename}{Daftar Gambar}}

%%%%% Mengganti label "Table" ke "Tabel" %%%%%
\addto\captionsenglish{\renewcommand{\tablename}{Tabel}}

%%%%% Mengganti label "List of Tables" ke "Daftar Table" %%%%%
\addto\captionsenglish{\renewcommand{\listtablename}{Daftar Tabel}}

\hypersetup{
	colorlinks=true, %set true if you want colored links
	linktoc=all,     %set to all if you want both sections and subsections linked
	linkcolor=blue,  %choose some color if you want links to stand out
	urlcolor=blue,   %url color
}

\geometry{
	a4paper,
	left=10mm,
	right=10mm,
	top=15mm,
	bottom=15mm,
}

\date{}

\hypersetup{citecolor=black}

\definecolor{LightGray}{gray}{0.95}

%\pagecolor[rgb]{0.1,0.1,0.1}
%\color[rgb]{1,1,1}

\lstset
{
	language=bash,
	breaklines=true,
	basicstyle=\tt\normalsize,
	frame = single
}

\begin{document}
	\frontmatter
	\begin{titlepage}
		\centering
		{\LARGE \bf Dokumentasi Kegiatan Komparasi Perangkat MiniDSP EARS dan Sound Level Meter}
		\vfill
		{\Large Achmadi ST MT}
		\vfill
		Update: {\today}
		\vfill
		\includegraphics[width=250pt]{images/logo/logoviblab}
		\vfill
		\vfill
		\vfill
	\end{titlepage}

	%%%%%%%%%%%%%%%%%%%%%%%%%%%%%%%%%%%%%%%%%%%%%%%%%%%%%%%%%%%%%%%%%

	\newpage
	\tableofcontents
	\listoffigures
	\listoftables

	%%%%%%%%%%%%%%%%%%%%%%%%%%%%%%%%%%%%%%%%%%%%%%%%%%%%%%%%%%%%%%%%%

	\mainmatter
	\newpage

	\chapter{Dokumentasi Kegiatan}

	Dokumen berikut menjelaskan kegiatan komparasi antara MiniDSP EARS dan Sound Level Meter.

	\section{Diagram}

	Berikut Diagram setup komparasi:

	\begin{figure}[!ht]
		\centering
		\includegraphics[width=0.5\textwidth]{images/diagram/tesearslm}
		\caption{Diagram Kegiatan}
	\end{figure}

	\newpage
	\section{Daftar Perangkat}

	Perangkat yang digunakan dalam kegiatan ini meliputi:

	\begin{enumerate}
		\item Sound Level Meter Rion NL-62
		\begin{figure}[!ht]
			\centering
			\includegraphics[width=0.5\textwidth]{images/perangkat/rion_nl62}
			\caption{Sound Level Meter}
		\end{figure}

		Website Resmi: \href{https://rion-sv.com/products/10005/NL620009}{Rion NL-62}

		\item Speaker BOSE Revolve Soundlink
		\begin{figure}[!ht]
			\centering
			\includegraphics[width=0.5\textwidth,angle=-90]{images/perangkat/bose}
			\caption{Speaker 360}
		\end{figure}

		Website Resmi: \href{https://www.bose.com/en_us/products/speakers/portable_speakers/soundlink-revolve-ii.html}{Bose Revolve Soundlink}

		\newpage
		\item MiniDSP EARS
		\begin{figure}[!ht]
			\centering
			\includegraphics[width=200pt,angle=-90]{images/perangkat/ears}
			\caption{MiniDSP EARS}
		\end{figure}

		Website Resmi: \href{https://www.minidsp.com/products/acoustic-measurement/ears-headphone-jig}{MiniDSP EARS}.
	\end{enumerate}

	Berikut setup keseluruhan dari kegiatan:

	\begin{figure}[!ht]
		\centering
		\includegraphics[width=0.7\textwidth]{images/perangkat/setup}
		\caption{Setup Keseluruhan}
	\end{figure}

	\section{Audio Analyzer}

	Untuk perangkat lunak Audio Analyzer, digunakan program \textit{Real-Time Analyzer}
	yang merupakan bagian dari paket software DSFF3 buatan Yoshimasha.
	Jika membutuhkan versi trial dapat didownload di \href{https://drive.google.com/drive/folders/1Z4cc4C_vb7BwU_MHDFyJTTVzQO7n7AHP?usp=share_link}{Audio\_Analyzer}.

	\begin{figure}[!ht]
		\centering
		\includegraphics[width=300pt]{images/dssf3/dssf3}
		\caption{DSSF3}
	\end{figure}

	\subsection{Kalibrasi EARS}

	Proses kalibrasi dilakukan pada software Yoshimasa dengan pengatuan level 76 dBA/80dB (sesuai dengan kondisi perekaman suara di dalam ruang diffuse).
	Rekaman ruang diffuse dapat didownload di \href{https://drive.google.com/drive/folders/1FfFU9WPwoisH7j-oDqF8G0ruZmY0lxM_?usp=share_link}{EARS\_Calibration}\\

	\textbf{Catatan:} Tersedia pula rekaman untuk Tone 125Hz, namun direkomendasikan menggunakan rekaman derau untuk mencakup semua frekuensi.\\

	Berikut Langkah Kalibrasi:
	\begin{enumerate}
		\item Sambung kabel USB EARS ke komputer.
		\begin{figure}[!ht]
			\centering
			\includegraphics[width=0.5\textwidth]{images/dssf3/laptop_ears}
			\caption{Kabel USB}
		\end{figure}

		\newpage
		\item Jalankan program \textbf{Realtime Analyzer}, kemudian buka tab fungsi \textbf{FFT Analyzer} dan \textbf{Recorder}.
		\begin{figure}[!ht]
			\centering
			\includegraphics[width=0.8\textwidth]{images/dssf3/fft_analyzer}
			\caption{FFT Analyzer}
		\end{figure}

		\item Pada Jendela Recorder, klik \textbf{Load} untuk memasukkan rekaman
		\begin{figure}[!ht]
			\centering
			\includegraphics[width=0.8\textwidth]{images/dssf3/load_file}
			\caption{Load File}
		\end{figure}

		\newpage
		\item Selanjutnya pada Jendela FFT Analyzer, Klik \textbf{Calibration}.
		Ikuti bagaimana panah dalam gambar di bawah dan atur nilai penyesuaian di 76dBA/80dB:

		\begin{figure}[!ht]
			\centering
			\includegraphics[width=\textwidth]{images/dssf3/adjustment}
			\caption{Adjustment}
		\end{figure}
	\end{enumerate}

	\newpage
	Antar-muka Real-Time Analyzer yang digunakan adalah \textit{FFT-Analyzer} pada tab \textit{Octave band}.

	\begin{figure}[!ht]
		\centering
		\includegraphics[width=\textwidth]{images/dssf3/fft}
		\caption{Octave band pada FFT analyzer}
	\end{figure}

	Pada tampilan Octave-Band, klik satu bar frekuensi untuk mengetahui nilai SPL (dB)
	yang terisolasi dengan frekuensi lain.

	\subsection{Noise Generation}

	Selain digunakan untuk mendapatkan nilai input, software DSSF3 juga dapat menghasilkan output Noise dan Tone-Generation.

	Berikut contoh langkahnya:

	\begin{enumerate}
		\item Sambungkan komputer/laptop ke Speaker (jika diperlukan).

		\item Jalankan program \textbf{Realtime Analyzer}, kemudian buka tab fungsi \textbf{Signal Generator}.

		\newpage
		\begin{figure}[!ht]
			\centering
			\includegraphics[width=0.7\textwidth]{images/dssf3/signalgen}
			\caption{Signal Generator}
		\end{figure}

		\item Dapat dipilih tab \textbf{Noise}, kemudian \textbf{Start}.
	\end{enumerate}

	\section{Hasil Komparasi}

	Selanjutnya dilakukan contoh uji komparasi untuk satu nilai white-noise,
	didapatkan SLM menunjukkan nilai 93.8 dBA,
	sedangkan dari EARS yang ditampilkan melalui software FFT Analyzer menunjukkan 94.66 dBA pada channel kiri (terdekat ke speaker).

	\begin{figure}[!ht]
		\centering
		\begin{subfigure}[t]{0.4\textwidth}
			\includegraphics[width=\textwidth]{images/hasil/spl_rta}
			\caption{Nilai EARS}
		\end{subfigure}
		\begin{subfigure}[t]{0.4\textwidth}
			\includegraphics[width=\textwidth]{images/hasil/spl_slm}
			\caption{Nilai SLM}
		\end{subfigure}

		\caption{Pembandingan EARS dan SLM}
	\end{figure}

	Dokumentasi berupa video singkat dapat dilihat di \href{https://drive.google.com/drive/folders/1SL2Z-jUqwHkQO1tEB37UERUyHxoY6Bd-?usp=share_link}{Google-Drive}
\end{document}
